\begin{established}[Short harmonic polynomials at the selected deep level]\label{hyp:deep-short-span}
Take the deep construction with \(N_*=g_0g_1 15^e\) and
\(\ell=L\). The nonterminal lift has positive signature blocks, \([E:\mathbb Q]/2>2\),
good integral models at every prime, and auxiliary level one.
The bounds on the level and the two fields are
\[
 \log N_*\le30\ell,\quad s_{\rm fld}(E)\le3\ell,\quad
 \log(fr)\le\ell,\quad s_{\rm fld}(F)\le391 Z_0^2 .
\]
For the family and the weights used in the short rules, the deep
weight bounds and \(Z_0\le e^{2\ell},\ \kappa\le e^{2\ell}\) give
\begin{equation}\label{eq:deep-short-weight-sizes}
 \begin{aligned}
 S(u_1)&\le Z_0^2(\kappa/T+1)\le e^{7\ell},\\
 S(v(j))&\le Z_0^2\kappa\le e^{6\ell}.
 \end{aligned}
\end{equation}
Both log sizes therefore satisfy the \(10\ell\) bound required by
the short-family estimate.
At \(u_2\) existence
of the auxiliary integral frame uses the ordinary full
lattice result: adding the remainder to \(u_1\) retains
positive consecutive gaps. The short frame entering the
multiplication estimate is at \(u_1\).

Write \(n=8(fr+1)\), \(d_a\) for the harmonic degree and
\[
 K_\theta=d_a+nfr/2.
\]
The even-column stable theta span
uses level classes generating the marked quotient, self-dual
auxiliary lattices \(\Lambda\) locally standard at the primes
of \(Q_S qd N_*\), and compatible auxiliary levels.
The polynomials are in the integral harmonic lattice of
the specified exponent array. They can be chosen with coefficient norm at most \(\exp(C_A Z_0^2 K_\theta)\). The auxiliary dimension satisfies \(n=8(fr+1)\ge8fr+4\);
the entries of the shifted positive block exceed \(n/2\).
Combining the shifts \(n_{0,\tau}\) in the theta weights with \eqref{eq:deep-short-weight-sizes}
and \(\log(fr)\le\ell\) gives the short-rule exponent bound
\[
 K_\theta=d_a+\tfrac12 nfr,\qquad \log K_\theta\le20\ell .
\]
\label{cor:deep-short-readings-1}
\end{established}
\begin{proof}
The level and coefficient-field estimates are those of Lemma~\ref{lem:deep-cover-field-size}; Proposition~\ref{prop:deep-cover-full-lattice} supplies the integral frame at \(u_2\). The powers of \(Z_0\) in \eqref{eq:deep-short-weight-sizes} allow both lists of weights, uniformly in the choice of \(j\). The stable span with a prescribed finite set of standard local lattices is constructed at \ref{cor:deep-cover-denominators-9}. We explain why its harmonic polynomials may be chosen short. Write \(D=d_a\) and \(r=p+q\). The constant polynomial suffices if \(D=0\).

Equip \(\Lambda\) with the real trace norm of its positive Hermitian form. Its trace-dual lattice is \(\mathfrak d_E^{-1}\Lambda\) by Hermitian self-duality. For \(d_E=|\operatorname{disc} E|\), multiplication by \(d_E\) sends this dual into \(\Lambda\), and every nonzero trace norm squared in \(\Lambda\) is a positive integer. Hence Euclidean transference in dimension \(2fn\) gives independent \(v_i\in\Lambda\) spanning over \(\mathbb R\) with
\[
        \lambda_1(\Lambda^*)\ge d_E^{-1},\qquad
        \max_i\|v_i\|\ll fn d_E.
\]
Here either convention for the complex trace norm changes only the absolute constant. \cite{banaszczyk}

For two block partitions \(a^+,a^-\), with trailing entries extended by zero, write the reference joint highest vector as
\[
  P_{a^+,a^-}=\prod_{s=1}^{p}\det(X_{[s],[s]})^{a_s^+-a_{s+1}^+}
          \prod_{t=1}^{q}\det(Y_{[t],[t]})^{a_t^--a_{t+1}^-},
\]
where the reference \(X,Y\) use orthogonal column frames. At each label the row-highest harmonics are generated by block products of leading minors evaluated on pairs of row frames \((A,C)\) with \(h(A,C)=0\). Indeed the mixed contractions of each such product vanish by orthogonality; at the coordinate frame the product is the highest vector of the required joint harmonic representation. Its column-group orbit spans that representation by highest-weight theory, with the determinant exponents equal to the successive differences of the prescribed rows, as in Section~\ref{subsec:theta-row-rank}. In particular, if \(V_a\) is the row-highest harmonic space, the column action satisfies
\[
      V_a=\operatorname{span}\{gP_{a^+,a^-}:g\in\mathrm{GL}_n\}
\]
at each complex label. The allowed coefficient embeddings occur in distinct conjugate pairs; after restriction of scalars their complexifications are independent factors.

Here is an integral parametrization of a dense set of these pairs. Choose row lists \(A,A'\) of length \(p\) and rows \(B_j\) of length \(q\), each a linear combination of the \(v_i\). Set \(H_{ik}=h(A'_k,A_i)\), \(b_{j,i}=h(B_j,A_i)\), taking the form linear in its first entry, and put
\[
 C_j=\det(H)B_j-\sum_k\bigl(\operatorname{adj}(H)b_j\bigr)_k A'_k .
\]
Hermitian integrality puts all these rows in \(\Lambda\) for integral coefficients. The vector of contractions is
\[
   \big(h(C_j,A_i)\big)_i=\det(H)b_j-H\operatorname{adj}(H)b_j=0 .
\] On the open set \(\det H\ne0\) this map onto the orthogonal complement is surjective in the variables \(B_j\), since its restriction there is scalar multiplication. The integral vectors \(v_i\) span the restriction of scalars over \(\mathbb Q\). After base change the parameters can therefore vary independently in each embedding, giving the dense set just used.

The entries of the resulting degree-\(D\) block product are polynomials of degree at most \(d=(2r+1)D\) in the integer parameters. Interpolate in one variable at a time on \(\{0,\ldots,d\}\). Indeed for every linear functional \(\ell_0\) on the coefficient space, successive one-dimensional interpolation gives
\[
   (\ell_0 P)|_{\{0,\ldots,d\}^{\,2fn(2p+q)}}=0
                  \quad\Longrightarrow\quad \ell_0P\equiv0 .
\]
Thus the grid spans.

With \(R=\max(1,2fn d\max_i\|v_i\|)\), the entries of a row in \(A,A',B\) have norm at most \(R\) at every complex embedding. Determinant expansion in the formula for \(C_j\) bounds each entry by \((p+1)!(nR^2)^p R\). Thus the coefficient \(\ell^1\) norm of each block product satisfies
\[
 \log\|P\|_1
 \ll D r\left(\log(2r n)+\log R\right)
 \ll D r\bigl(f\ell+\log(1+D)\bigr)
 \ll Z_0^2 K_\theta.
\]
The second estimate uses \(\log d_E\le 7f\ell\), and the last uses \(Z_0=fr\ell\) and the weight bounds above. The same calculation applies at every embedding of the coefficient field.
\end{proof}

\begin{proposition}[An integral short span on the deep cover]
\label{prop:deep-short-integral-span}
For the deep cover with the data of Lemma~\ref{hyp:deep-short-span},
each even-weight spanning rule \(y\) has an integral
multiple \(sZ_*^2y\), with
\[
     \log\|s Z_*^2 y\|_{\mathrm P}
       \le s_{\rm fld}(F)+O(Z_0^2K_\theta).
\] The determinant-ideal scalar may depend
on the auxiliary lattice. The second scalar depends only on the
level and harmonic degree. The flat-line shift preserves the
rational span with nonnegative order at every integral frame and the
additional archimedean cost in \eqref{eq:parallel-return-norm}.
\end{proposition}

\begin{proof}
\proofstep{Descent of the even-weight span.}
Use the construction at \ref{cor:deep-cover-denominators-9} with
the auxiliary integer divisible by \(Q_S qd\). It realizes the
short span on the identity cover with the standard auxiliary
lattice at every prime dividing \(Q_S qd N_*\). Expand the varying polynomial in short integral harmonic polynomials for that lattice. The resulting theta section depends linearly on the polynomial. The finite class indicator gives the
marked Fourier coefficients, and the orbit of the component
under coefficient embeddings has the prescribed uniformized
Hermitian space. These conditions allow the marked sections and the span of their short theta rules to be compared in the same coefficient space over \(F\).
To state the descent explicitly, let \(V\) be that space and
\(W\subset V\) the span of the specified short rules. If a section \(y\in V\) has complexification in \(W\otimes_F\mathbb C\),
then
\[
 (y+W)\otimes1=0\ \text{in }(V/W)\otimes_F\mathbb C
                 \quad\Longrightarrow\quad y\in W .
\]

\proofstep{The determinant-ideal scalar.}
Two different scalars enter. One depends on \(\Lambda\).
Let \(I\) be the index of its intersection with the
standard lattice in \(\Lambda\), and \(\mathfrak a\) the
index ideal, of norm \(I\). In each embedded determinant
line the intersection contributes \(\tau\mathfrak a\).
Take \(B=\prod_\tau(\tau\mathfrak a\,\mathcal O_F)^{m_\tau}\)
with the signature multiplicities. Its norm is
\(I^{r[F:\mathbb Q]/2}\) since each conjugate pair
has total multiplicity \(r\).

In the totally complex
coefficient field, apply the convex body theorem to
\(B^{-1}\) with disk radii
\(\exp(s_{\rm fld}(F)) I^{-r/2}\). The lattice
covolume is \(2^{-[F:\mathbb Q]/2}
\sqrt{|\operatorname{disc} F|}/\mathrm N B\).
The disk product exceeds \(2^{[F:\mathbb Q]}\) times this
covolume, by the field-size bound and \(\pi>2\).
There results \(s\ne0\) in \(B^{-1}\) with
\[
   |\sigma s|^2 I^r\le \exp(2s_{\rm fld}(F))
           \qquad\text{for every }\sigma.
\]
Membership is precisely the determinant-ideal inclusion
used for the integral short rule. The same \(s\) pays
for the factor \(I^{r/2}\) in its analytic bound.

\proofstep{A second scalar for finite prime and level denominators.}
The other scalar can be fixed before the auxiliary lattice
and the class tuple are chosen. Set
\[
    Z_* = N_*^{frn+d_a}\,|\operatorname{disc}E|^{d_a}.
\]
The level class average is cleared by its cardinality
\(N_*^{frn}\), and the polynomial and eigenframe factors
by
\[
 N_*^{d_a}|\operatorname{disc}E|^{d_a}.
\]

In the finite order test write the raw rule \(y\) as
\(\beta_\Lambda\) times the short rule in determinant
coordinates. The class indicator at each label has integral
Fourier coefficient (a root of unity). With \(s\) in the inverse index ideal, the coprime torsion comparison gives 
\[
 v(sZ_*y)\ge0\qquad(v\nmid N_*).
\]
 At a prime dividing \(N_*\) the auxiliary
lattice is locally standard, so \(s\) is integral there.
Use the chart comparison at that prime for the short
rule in determinant coordinates, and multiply back by
\(\beta_\Lambda\).

For a prime dividing the level once
the class-sum bound clears this with \(Z_*\); for one
whose square divides the level the trivial-lattice
comparison clears it with \(Z_*^2\). In the latter comparison the class bound is applied after a finite extension. For an element \(a\) of the base and integer-normalized orders upstairs and downstairs, the normalization is
\[
 v_{\rm base}(a)=e^{-1}v_{\rm ext}(a),
\]
where \(e\) is the ramification index. Summing integral Fourier coefficients
preserves the lower order and descends it to the original
test. Combining the local tests gives
\[
 v(sZ_*^2y)\ge0\qquad(\text{every finite test place}),
\]
including the two primes of the growing principal depth. Moreover
\[
  \log Z_*^2
    =2(frn+d_a)\log N_*+2d_a\log|\operatorname{disc}E|
    \ll Z_0^2 K_\theta
\]
since \(frn\le2K_\theta\) and
\(\log|\operatorname{disc}E|\le7f\ell\).
\label{cor:deep-short-readings-2}

\proofstep{The common complex norm bound.}
For the complex comparison choose, at an embedding with
positive supremum, an orthonormal flag frame reaching
at least half the supremum and with squared first minimum
at least \(e^{-80\ell}\). This is the level point
estimate for spread at most \(3K_\theta\): its level
cutoff is at least \(e^{-57\ell}\) and the loss in
squared first minimum is at most \(\operatorname{rd}(E)^2\).
The Fourier transform of the point indicator has absolute
value one at each label.

On that frame the classical
self-dual theta norm estimate with this first minimum,
the polynomial coefficient bound and its determinant factor
gives
\[
   \|sy\|_{\sup}\le
   \exp(s_{\rm fld}(F))\,
   \exp\big((C_A+C_B)Z_0^2 K_\theta\big).
\]
The common Lie factor in this estimate is
\[
 (2\pi\xi_R)^{-d_a/2};
\]
the index factor is paid
by the preceding inequality for \(s\). The uniformization
uses the probability mass, so the same bound holds for
the Petersson norm on each component. Multiplication by
\(Z_*^2\) and \(K_\theta\ge1\) absorb the remaining
field and integer costs. All these scalar multiples are
nonzero multiples of the respective spanning rules.
Restriction to \(Z\) therefore gives an integral
Petersson-bounded span at the even theta weight.
\label{cor:deep-short-readings-3}

\paragraph{Making the flat weight return integral.}
The key unit and the parallel return are used with
different estimates. For the parallel return let
\(r_\sigma>0\) be its constant complex norms and
\(\rho_w\) its finite orders before scalar adjustment.
They have finite order support, conjugate embeddings
have the same norm, and outside the bad rational primes
the orders are integers by the unramified frame test.
In \(F\) form the fractional ideal
\[
      J=(b_{\rm bad})\prod_w\mathfrak p_w^{\lfloor-\rho_w\rfloor},
             \qquad b_{\rm bad}=\prod_{l\,\mathrm{bad}} l .
\]
Every nonzero \(a\in J\) satisfies
\(v_w(a)+\rho_w\ge0\). Off the bad primes this follows
by integrality; at a bad prime \(v_w(b_{\rm bad})\ge1\)
covers the rounding.

If \(h=[F:\mathbb Q]\), then
\[
 \tfrac1h\log\mathrm N J\le\log b_{\rm bad}
            -\tfrac1h\sum_w\rho_w\log\mathrm N w.
\]
Use a symmetric Minkowski body with radii
\(A/r_\sigma\), where
\(A=(5/2)(\mathrm N J\sqrt{|\operatorname{disc}F|}
\prod_\sigma r_\sigma)^{1/h}\). Its volume exceeds
the Minkowski threshold for \(J\). The resulting nonzero
element gives a single scalar at all embeddings. Taking
logs yields
\begin{equation}\label{eq:parallel-return-norm}
  \log(|\sigma a|r_\sigma)\le
  -\deg_{\rm flat}+\tfrac12\log\operatorname{rd}(F)
             +\log b_{\rm bad}+1 .
\end{equation}
Here the degree uses the archimedean negative log norms
and the finite additive orders. The preceding flat
calculation bounds the negative degree by the expression
linear in the absolute exponents below. Thus the
parallel twist has nonnegative order at every integral
frame and a simultaneous upper norm with this loss. It
is invertible on the generic component, so multiplication
preserves the span. \label{cor:deep-short-readings-4}
\end{proof}
